% CREATED BY DAVID FRISK, 2016
\chapter{Discussion}

The superior reconstruction correlation achieved by UniEdgeCodec (0.960 $\pm$ 0.007 on EMOB) compared to its competitors can be understood through its design objective of optimizing for perceptual quality, particularly for irregular and discontinuous signals. Since perceptual quality is fundamentally correlated with reconstruction correlation, as can be seen by the visualizations of the reconstructed signals, the explicit optimization for this objective naturally leads to improved correlation metrics across all signal types. This consistent preservation of temporal structure is especially beneficial for downstream anomaly detection tasks, where classification reliability depends on the accurate capture of temporal patterns. The algorithm's architectural advantage, preventing signals from competing within the same codebook, combined with the smoothing effect of the loss function, contributes to the training stability observed across all random seeds. These factors synergistically produce both the lowest metric variance and the most reliable convergence behavior across the experimental runs.

The mixed and seed-dependent results of TAEdgeCodec reveal fundamental challenges in multi-objective neural compression. The more complex optimization landscape that arises from jointly learning compression and downstream task performance creates inherent trade-offs: the model must allocate limited representational capacity between reconstruction fidelity and task-specific features. This explains the observed pattern of worse reconstruction (0.765 $\pm$ 0.035 correlation on EMOB) compared to UniEdgeCodec, instability in utility retention (67.0 $\pm$ 13.8\% on EMOB), and occasional convergence failures such as the early collapse observed in seed 100 for the Volvo EMOB data. The high utility retention sometimes achieved on the appliances energy dataset (89.4 $\pm$ 0.5\%) suggests that when optimization objectives align well with the downstream task, the approach can succeed. However, the substantially higher variance across seeds indicates that the balance between reconstruction and task optimization is fundamentally unstable without the additional architectural constraints provided by UniEdgeCodec's per-signal codebook design.

EdgeCodec's strong reconstruction performance across both datasets (0.780 $\pm$ 0.034 on EMOB, 0.635 $\pm$ 0.061 on appliances energy) is consistent with its primary optimization objective of reconstruction fidelity. However, its utility retention results reveal an important signal-dependent phenomenon. For the EMOB classification task, the model faces a trade-off between reconstructing smooth signals like temperature and irregular signals like current (``Amps''), which tend to share compressed representations in the shared codebook. While EdgeCodec reconstructs temperature adequately, supporting anomaly detection on channels where temperature correlates with anomalies, this comes at the cost of poor reconstruction on channels with sparse, irregular patterns like current fluctuations. UniEdgeCodec mitigates this problem by dedicating separate codebooks to each signal, allowing simultaneous optimization for both signal types. The intermediate training stability of EdgeCodec (between the stable UniEdgeCodec and unstable TAEdgeCodec) suggests that the simplified, single-codebook optimization landscape is easier to navigate than TAEdgeCodec's multi-objective problem, but lacks the architectural advantages that lead to UniEdgeCodec's superior convergence properties.

The TCN-RNN baseline demonstrates that continuous latent representations become increasingly problematic at high compression ratios. Operating at only 3.0x compression with a correlation of 0.726 on EMOB—substantially below the 200x+ compression of discrete methods—reveals that continuous bottleneck compression is inherently limited for sequence compression tasks. While discrete latent representations provide a clear advantage, the baseline results also underscore that model objectives matter significantly. The discrete methods' superior performance stems not only from discrete latents, but from the specific optimization objectives designed for each architecture.

The relationship between reconstruction quality and downstream utility retention exhibits important task-dependent characteristics. For the EMOB anomaly detection task, reconstruction metrics serve as reliable predictors of utility retention because the task fundamentally depends on capturing sparse, discontinuous events in the signal. When these critical events are not reliably reconstructed—as evidenced by poor MSE—downstream anomaly detection inherently suffers. This explains the correlation visible in Table~\ref{tab:uni-ta-per-seed} between reconstruction metrics and utility retention on EMOB. In contrast, the appliances energy regression task is more robust to compression artifacts because it primarily depends on capturing overall signal trends rather than precise reconstruction of individual spikes or discontinuities. This task-dependent behavior highlights that compression suitability cannot be universally assessed through reconstruction metrics alone; rather, the choice of compression method must be informed by the specific requirements of the downstream application.

% Actionable Suggestion
